\chapter{Units and Conversions}

In 1999, NASA lost the \$125 million Mars Climate Orbiter because one team used pounds and another used newtons. Units matter.

You are already juggling a lot of them: grams for mass, newtons for force, meters for distance, seconds for time, joules for energy. Worse, each kind of measurement has several units. Distance alone comes in inches, feet, miles, kilometers, and light-years. This chapter shows you how to move between them without getting lost.

\begin{mdframed}[style=important, frametitle={Some Equivalencies}]

\begin{tabular}{r | l}
  \hline
  \multicolumn{2}{c}{\textbf{Distance}}\\
  1 mile & 1.6093 kilometers \\
  1 foot & 0.3048 meters \\
  1 inch & 2.54 centimeters \\
  1 light-year & $9.461 \times 10^{12}$ kilometers\\
  \hline
  \multicolumn{2}{c}{\textbf{Volume}}\\
  1 milliliter & 1 cubic centimeter \\
  1 quart & 0.9464 liters \\
  1 gallon & 3.7854 liters \\
  1 fluid ounce & 29.6 milliliters \\
  \hline
  \multicolumn{2}{c}{\textbf{Mass}}\\
  1 pound & 0.4535924 kilograms\\
  1 ounce & 28.35 grams\\
  1 metric ton & 1000 kilograms \\
  \hline
  \multicolumn{2}{c}{\textbf{Force}}\\
  1 newton & 1 kilogram meter per sec$^2$\\
  \hline
  \multicolumn{2}{c}{\textbf{Pressure}}\\
  1 pascal\index{pascal} & 1 newton per square meter \\
  1 bar & 0.98692 atmosphere \\
  1 pound per square inch & 6895 pascals \\
  \hline
  \multicolumn{2}{c}{\textbf{Energy}}\\
  1 joule & 1 newton meter \\
  1 calorie & 4.184 joules \\
  1 kilowatt-hour\index{kilowatt-hour} & $3.6 \times 10^{6}$ joules  \\
\end{tabular}\index{units table}

(You do not need to memorize these! Just remember that this page is here.)
% Suggest putting this in the front or back of the book, maybe create a print out

\end{mdframed}

In the metric system, prefixes are often used to express a multiple. Here are the common prefixes:\index{metric system!prefixes}

\begin{mdframed}[style=important, frametitle={Common Prefixes for Metric Units}]

\begin{tabular}{r | l}
giga  & $\times 10^{9}$\\
mega  & $\times 10^{6}$\\
kilo  & $\times 10^{3}$\\
milli  & $\div 10^{3}$\\
micro  & $\div 10^{6}$\\
nano  & $\div 10^{9}$\\
\end{tabular}

(These are worth memorizing. Here is a mnemonic for the most common ones: ``King Henry Died By Drinking Chocolate Milk.'' The letters stand for kilo ($\times 10^{3}$), hecto ($\times 10^{2}$), deca ($\times 10^{1}$), base unit such as the gram ($\times 10^{0}$), deci ($\times 10^{-1}$), centi ($\times 10^{-2}$), and milli ($\times 10^{-3}$).)
\end{mdframed}

\section{Conversion Factors}

Here is the one trick you need.\index{conversion factors}

Suppose someone asks, ``How many miles are in 0.23 light-years?'' From the table, 1 mile is 1.6093 kilometers and 1 light-year is $9.461 \times 10^{12}$ kilometers. How do you get from one to the other?

Turn each equality into a fraction that equals 1:

$$\frac{1 \text{ mile}}{1.6093 \text{ km}} = \frac{1.6093 \text{ km}}{1 \text{ mile}} = 1$$

and

$$\frac{1 \text{ light-years}}{9.461 \times 10^{12} \text{ km}} = \frac{9.461 \times 10^{12} \text{ km}}{1 \text{ light-years}} = 1$$

We call these fractions \textit{conversion factors}.

Now, your problem is

$$0.23 \text{ light-years} \times \textit{ Some conversion factors} = ? \text{ miles}$$

Note that when you multiply fractions together, things in the numerators can cancel with things in the denominator:

$$\left( \frac{31\pi}{47} \right) \left( \frac{11}{37\pi}\right) = \left(\frac{31\cancel{\pi}}{47}\right) \left( \frac{11}{37\cancel{\pi}}\right) = \left(\frac{31}{47} \right) \left( \frac{11}{37} \right)$$

Units cancel the same way\index{unit cancellation}. Multiplying by 1 never changes a value, so you can multiply by as many conversion factors as you like. Choose them so the units you do not want cancel out:

\textbf{Worked example: converting 0.23 light-years to miles.}
\begin{multline*}
  0.23 \text{ light-years} \left( \frac{9.461 \times 10^{12} \text{ km}}{1 \text{ light-years}} \right) \left( \frac{1 \text{ mile}}{1.6093 \text{ km}} \right) = \\
  0.23 \text{ \cancel{light-years}} \left( \frac{9.461 \times 10^{12} \text{ \cancel{km}}}{1 \text{ \cancel{light-years}}} \right) \left( \frac{1 \text{ mile}}{1.6093 \text{ \cancel{km}}}\right) = \frac{(0.23)(9.461 \times 10^{12})}{1.6093} \text{ miles} \approx 1.35 \times 10^{12} \text{ miles}
\end{multline*}

\begin{Exercise}[title={Simple Conversion Factors}, label=simple_conversion_factors]

  How many calories are in 4.5 kilowatt-hours?

\end{Exercise}
\begin{Answer}[ref=simple_conversion_factors]

  $$4.5 \text{ \cancel{kWh}} \left( \frac{3.6 \times 10^{6} \text{ \cancel{joules}}}{1 \text{ \cancel{kWh}}} \right) \left( \frac{1 \text{ calories}}{4.184 \text{ \cancel{joules}}}\right) = \frac{(4.5)(3.6 \times 10^6)}{4.184} = 3.87 \times 10^6 \text{ calories}$$

\end{Answer}

\section{Conversion Factors and Ratios}

Conversion factors also work on ratios. For example, if you are told
that a bug is moving 0.5 feet every 120 milliseconds, what is that in
meters per second?

The problem then is

$$\frac{0.5 \text{ feet}}{120 \text{ milliseconds}} = \frac{\text{? m}}{second}$$

So you will need conversion factors to replace the ``feet'' with ``meters'' and to replace ``milliseconds'' with ``seconds'':

\begin{multline*}
\left(\frac{0.5 \text{ \cancel{feet}}}{120 \text{ \cancel{milliseconds}}}\right) \left( \frac{0.3048 \text{ meters}}{1 \text{ \cancel{feet}}} \right) \left( \frac{ 1000 \text{ \cancel{milliseconds}}} {1 \text{ second}}\right) = \frac{(0.5)(0.3048)(1000)}{120}\text{ m/second} \approx 1.27 \text{ m/second}
\end{multline*}

\begin{Exercise}[title={Conversion Factors}, label=conversion_factors]

The hole in the bottom of the boat lets in 0.1 gallons every 2 minutes. How many milliliters per second is that?

\end{Exercise}
\begin{Answer}[ref=conversion_factors]

  \begin{multline*}
    \frac{0.1 \text{ \cancel{gallons}}}{2 \text{ \cancel{minutes}}}
  \left( \frac{3.7854 \text{ \cancel{liters}}}{1 \text{ \cancel{gallons}}} \right)
  \left( \frac{1000 \text{ milliliters}}{1\text{ \cancel{liters}}}\right)
  \left( \frac{1 \text{ \cancel{minutes}}}{60 \text{ seconds}} \right) = \\
  \frac{(0.1)(3.7854)(1000)}{(2)(60)} \text{ ml/second} = 3.1545 \text{ ml/second}
  \end{multline*}

\end{Answer}

\section{When Conversion Factors Fail}

Conversion factors only work when the units being converted are
proportional to each other. Gallons and liters\index{proportional units}, for example, are
proportional to each other: If you have $n$ gallons, you have $n
\times 3.7854$ liters.

\index{Celsius}\index{Fahrenheit}Degrees Celsius and degrees Fahrenheit are \textit{not} proportional. Zero degrees Celsius is 32 degrees Fahrenheit, not zero. If your soup is $n$ degrees Celsius, it is $n \times
\frac{9}{5} + 32$ degrees Fahrenheit. That ``$+ 32$'' is why no conversion factor can do this job.

\section{SI units}
Scientists everywhere use one standard set of units, so that a measurement made in Tokyo means the same thing in Toronto. It is called the \textbf{International System of Units}, or \textbf{SI}\index{SI units}. It has seven base units:

\begin{tabular}{r | l}
Time  & second ($s$) \\
Length  & meter ($m$)\\
Mass  & kilogram ($kg$)\\
Electric Current  & ampere ($A$)\\
Temperature  & Kelvin\index{Kelvin} ($K$)\\
Amount of substance  &  mole ($mol$)\\
Luminous intensity  & candela ($cd$)\\
\end{tabular}

Every other unit\index{derived units} is built from these seven.
% ADD (maybe?) the constants behind these 7 units... feels to complicated https://en.wikipedia.org/wiki/International_System_of_Units#SI_defining_constants


You will often see units written with exponents. Velocity is meters \textit{per} second, which you can write as a fraction or with seconds to a negative power:
$$v = \dfrac{m}{s} = m \cdot s^{-1}$$

Acceleration, meters per second per second, works the same way:
$$a = \dfrac{m}{s^2} = m \cdot s^{-2}$$

\subsection{Converting from SI units}
Many everyday units are SI units in disguise, built from the base units:

\begin{table}[!ht]
    \centering
    \begin{tabular}{|c|c|c|c|}
    \hline
        Physical Quantity & Name & Symbol & Expression in Base Units \\ \hline
        Force & Newton & N & $kg\cdot m\cdot s^{-2}$ \\ \hline
        Energy/Work & Joule & J & $kg\cdot m^2\cdot s^{-2}$ \\ \hline
        Power & Watt & W & $kg\cdot m^2\cdot s^{-3}$ \\ \hline
        Frequency & Hertz & Hz & $s^{-1}$ \\ \hline
        Electric Charge & Coulomb & C & $A\cdot s$ \\ \hline
    \end{tabular}
\end{table}

You can derive all of these yourself. Try the newton: force is mass times acceleration, so what are its base units?

\begin{Exercise}[title={Miles to Feet}, label=feet_to_miles]
How many feet are in 3 miles? Use 1 mile = 1.6093 kilometers and 1 foot = 0.3048 meters.
\end{Exercise}
\begin{Answer}[ref=feet_to_miles]
Convert miles to meters: $3 \text{ mile} \times \frac{1.6093 \times 10^{3} \text{ m}}{1 \text{ mile}} = 4827.9$ m. Then
convert meters to feet: $4827.9 \text{ m} \times \frac{1 \text{ ft}}{0.3048 \text{ m}} \approx 15{,}840$ ft.
\end{Answer}
